<<<P000775>%
Предполагается, что каждый <<<P000855>{typeName}, используемый в типе, упомянутом в этом разделе,
Не имеет ошибки компиляции и обозначает тип.

<<<P000856>{%>
То есть, никакие отношения субтипирования не могут быть доказаны.
Тип, который содержит или содержит неопределенное имя
или имя, которое обозначает нечто иное, чем тип.
Обратите внимание, что для определения отношений субтипирования нет необходимости.
Каждый тип удовлетворяет заявленным границам.
Отношение субтипирования не зависит от границ.
Однако, если предпринимается попытка доказать связь подтипа
и один или более <<<P000857>{typeName}s получает фактический список аргументов типа
длина которого не соответствует декларации
(включая случаи, когда приводятся аргументы определенного типа)
неродовой класс,
и случай, когда возникает общий класс,
Но никаких аргументов не приводится.
Тогда попытка доказать отношения просто проваливается.
?

<<P000776>%
Правила на рисунке ~<<<P000513> использовать
Символ <<P000550> Чтобы обозначить данное знание о
границы типовых переменных.
<<P000004> Это частичная функция, которая отображает переменные типа на типы.
В заданном месте, где переменные типа в области применения
<<<P000858><
<<P000859>{(как указано в прилагаемых классах и/или функциях)},
Мы определяем окружающую среду следующим образом:
<<P000005>.
<<<P000860>{%>
То есть, <<<P000006> и так далее,
<<<P000007> не определен при применении к переменной типа <<<P000008>
которого нет в <<<P000009>.%
?
Когда правила используются, чтобы показать, что существует определенный подтип отношений,
Это начальное значение <<<P000010>.

<<P000777>%
Если встречается общий тип функции, расширение <<<P000011> используется,
Как показано в правилах ~\SrnPositionalFunctionType
и ~<<<P000862>
Изображение: P000514.
Расширение сред использует оператор <<<P000551>,
который является оператором, производящим объединение разрозненных множеств,
и отдает приоритет правой руке операнда в случае конфликтов.

<<<P000863><%
Так
$<<<P001475> <<P000566> <<P000864> <<P000567>, <<P000568> <<P000865> <<P000569> <<<P001476> <<P000866>
<<P001477> <<P000570> <<P000867> <<P000571> <<<P001478>
<<<P001479> <<P000572> <<P000868> <<P000573>, <<P000574> <<P000869> <<P000575>, <<P000576> <<P000870> <<P000577> <<<P001480>
и
$<<<P001481> <<P000578> <<P000871> <<P000579>, <<P000580><<P000872> <<P000873> FutureOr<List<double>{}> <<P001482> <<P000874>
<<P001483> <<P000581> <<P000875><<P000582> <<<P001484>=
<<P001485> <<P000583> <<P000876> <<P000584>, <<P000585> <<P000877> <<P000586> <<<P001486>$.
Обратите внимание на оператора <<<P000013> Это связано с масштабами и тенью,
без связи, например, с подтипами или методом экземпляра overriding.%
?

<<P000778>%
В данной спецификации мы часто ссылаемся на
отношения подтипа и присваиваемость
Не говоря уже об окружающей среде,
<<<P000878><<<P000879> {T}}
Это делается только тогда, когда конкретное место в коде находится в фокусе.
Окружающая среда – это то, что получается.
путем отображения каждой переменной типа в объеме в этом месте
Объявлено его обязательным.

<<<P000779>%
Каждое правило на рисунке <<<P000515> имеет горизонтальную линию,
Слева от которого находится <<P000552> указывается;
Под горизонтальной линией находится суждение, которое
<<<P000880>{заключение}{правило!заключение}
Правило,
над горизонтальной линией находится ноль или более
<<<P000881>{предпосылки}{правило!предпосылка}
Правило,
которые, как правило, также являются подтипом суждения.
Если это не относится к определенной предпосылке
Мы четко указываем значение.

<<<P000882>{%>
Обоснование вышеупомянутого правила,
обозначает последовательную замену метавариабельных
посредством фактических синтаксических терминов, обозначающих типы повсюду в правиле,
То есть, как в помещениях, так и в заключении, одновременно.
?

<<P000883> Быть подтипом
<<P000836>

<<<P000780>%
Тип <<<P000014> Показано, что <<<P000553> другого типа <<<P000015>
в окружающей среде <<<P000016> путем предоставления
Инстанциация правила <<<P000017> чей вывод является
<<P000884><<P000885> S}{T}}{<<P000018>>>@\SubtypeStd S}{T}}
наряду с инсталляциями правил, показывающими
каждого из помещений <<<P000019>,
Продолжается до тех пор, пока не будет достигнуто правило, не имеющее помещения.

<<<P000887>{%>
Для правила <<<P000888> обратите внимание, что <<<P000587> тип
является подтипом всех не-<<<P000020> типы,
Хотя на самом деле он не расширяет и не реализует эти типы.
Другие типы эффективно рассматриваются так, как если бы они были недействительными.
Что делает нулевой объект (<<P000516>) Присваивается им. %
?

<<<P000781>%
Первая предпосылка в
Правила ~\SrnLeftTypeAlias и ~<<<P000890>
Это своеобразная декларация.
Эта предпосылка удовлетворяется в каждой из следующих ситуаций:

<<P000483>
<<P000891> Объявляется псевдоним недженерического типа под названием <<<P000021>.
В данном случае <<<P000022> равна нулю,
Нет никаких предположений о существовании
любых формальных типовых параметров,
Списки аргументов фактического типа опускаются везде в правиле.
<<P000892> Мы можем выбрать <<<P000023> и <<P000893>{X}{1}{s} таким образом, чтобы следующее:
Общий псевдоним типа <<<P000024> объявляется,
Формальные параметры типа <<P000894> X}{1}{s}
<<<P000895>{%>
Каждый формальный параметр типа <<<P000025> Может быть, у него есть связь,
Но границы никогда не используются в этом контексте.
Поэтому мы не вводим метавариабельности для них.%
?
<<P000498>

<<P000782>%
Правило ~\SrnRightFunction В качестве предпосылки, что <<<P000026> является функциональным типом.
Это означает, что <<<P000027> является одной из форм, введенных в
Раздел <<<P000517>.
<<<P000897>{%>
Это то же самое, что и формы, которые встречаются на высшем уровне.
в выводах из
Правило ~\SrnPositionalFunctionType{}
Правило ~\SrnNamedFunctionType%
?

<<P000783>%
В правилах ~\SrnCovariance и ~<<<P000901>,
Первое условие – это классовая декларация.
Эта предпосылка удовлетворяется в каждой из следующих ситуаций:

<<P000484>
<<P000902> Объявляется необычный класс с именем <<<P000028>.
В данном случае <<<P000029> равна нулю,
Нет никаких предположений о существовании
любых формальных типовых параметров,
Списки аргументов фактического типа опускаются везде в правиле.
<<P000903> Мы можем выбрать <<<P000030> и <<P000904>{X}{1}{s} таким образом, чтобы следующее:
Общий класс, названный <<<P000031> объявляется,
Формальные параметры типа <<P000905> X}{1}{s}
<<<P000906><<<<<<<<<<<<<<<<<<<<<<<<<<<<<<<<<<<<<<<<<<<<<<<<<<<<<<<<<<<<<<<<<<<<<<<<<<<<<<<<<<<<<<<<<<<<<<<<<<<<<<<<<<<<<<<<<<<<<<<<<<<<<<<<<<<<<<<<<<<<<<<<<<<<<<<<<<<<<<<<<<<<<<<<<<<<<<<<<<<<<<<<<<<<<<<<<<<<<<<<<<<<<<<<<<<<<<<<<<<<<<<<<<<<<<<<<<<<<<<<<<<<<<<<
Каждый формальный параметр типа <<<P000032> Может быть, у него есть связь,
Но границы никогда не используются в этом контексте.
Поэтому мы не вводим метавариабельности для них.%
?
<<P000499>

<<P000784>%
Вторая предпосылка правила ~<<<P000907> указывает, что
параметризированный тип \code{<<P000033>>><\ldots{}>> принадлежит
\IndexCustom{<<P000911>>>{C}}{суперинтерфейсы(C)@<<P000912>>>{ С.
Семантическая функция \Superinterfaces{\_} применяется к общему классу <<<P000034> доходность
Набор прямых суперинтерфейсов <<<P000035>
(P000518).

<<<P000914>{%>
Обратите внимание, что один из прямых суперинтерфейсов <<<P000036> это
интерфейс суперкласса <<<P000037>,
Это может быть приложение для смешивания
(<<P000519>),
в каком случае <<P000038> В правиле есть
синтетический класс, определяющий
Семантика этого смешанного приложения.%
?

<<<P000915>{%>
Последняя предпосылка правила ~<<<P000916>
Подменяет аргументы фактического типа \List{S}{1}{s}
Формальные параметры типа \List{X}{1}{}
<<<P000919>> T}{1}{m} может содержать формальные параметры типа.%
?

<<<P000920>{%>
Правила <<<P000921> и ~\SrnSuperinterface
применимы к интерфейсам,
Также их можно использовать с классами.Потому что негенерический класс <<P000039> который используется как тип
обозначает интерфейс <<<P000040>,
Для параметризированного типа
<<<P000923><<<P000041><<<<P000924> T}{1}{k}>
где <<<P000042> обозначает общий класс.%
?


<<P000925> Неофициальные описания правил подтипа
<<P000837>

<<<P000926>{%>
В этом разделе дается неформальное и ненормативное описание естественного языка.
Каждое правило на рисунке <<<P000520>.

В описаниях используются номера правил, чтобы сделать соединение явным,
и добавляет имена к правилам, которые могут быть полезны для понимания
Роль, которую играет каждое правило.

В следующем, многие правила содержат мета-вариабельность.
(<<P000521>)
как <<<P000043> и <<P000044>,
И всегда можно говорить о произвольных типах.
Например, правило ~\SrnRightFutureOrA{} Говорит, что
Тип <<<P000045> <<<P000928><<<P000588> <<<P000929>>,
Это означает, что для любых произвольных типов <<P000047> и <<P000048>,
Показать <<<P000049> является подтипом <<<P000050> Достаточно показать, что <<P000051> это
Подтип <<<P000589>.

Другой пример — формулировка в правиле ~\SrnReflexivity{}:
\ldots{} в любой среде <<<P000053>,
что указывает на то, что правило может применяться независимо от того, какие привязки
Переменные типа к границам существуют в окружающей среде.
Следует отметить, что окружающая среда имеет значение даже при соблюдении правил.
где оно просто заявлено как простое <<P000054> в заключение
в одном или нескольких помещениях,
Поскольку доказательство этих предпосылок может прямо или косвенно
Включает применение правила, в котором используется окружающая среда.

<<P000932>>><<P000933>>>#1#2{<<P000934>>>[#1]{<<P000935>>>{#2:}}}
<<P000485>
\Item{<<<P000937>> Рефлексивность
Каждый тип является подтипом самого себя, в любой среде <<P000055>.
В следующих правилах, кроме нескольких,
Это правило действует в любой среде.
Окружающая среда никогда не используется открыто.
Поэтому мы не будем повторять этого.
\Item{\SrnTop}} Вершина!
Каждый тип является подтипом <<<P000590>,
Каждый тип является подтипом <<<P000940>,
Каждый тип является подтипом <<<P000941>.
Это означает, что эти типы эквивалентны
в соответствии с отношением подтипа.
Мы обозначаем эти типы,
и другие с таким же свойством (например, <<P000591>),
как лучшие типы
(<<P000522>).
\Item{<<P000943>>>}
Каждый тип является супертипом P000056.
\Item{<<P000945>>>}
Все типы, кроме <<<P000057> Это супертип P000592.
\Item{\SrnLeftTypeAlias> Тип Alias Left
Применение псевдонима типа к некоторым фактическим аргументам типа
Подтип другого типа <<P000058>
если расширение псевдонима типа на тип, который он обозначает
Это подтип <<<P000059>.
Обратите внимание, что псевдоним необщего типа обрабатывается с помощью разрешения <<<P000060>.
\Item{\SrnRightTypeAlias}} Тип Alias Right
Тип <<<P000061> является подтипом приложения псевдонима типа
Если <<P000062> является подтипом
расширение псевдонима типа на тип, который он обозначает.
Обратите внимание, что псевдоним необщего типа обрабатывается с помощью разрешения <<<P000063>.
\Item{\SrnLeftFutureOr}} Левое будущее или
Тип <<<P000593> является подтипом данного типа <<P000065>
Если <<P000066> является подтипом <<<P000067> и <<<P000594> является подтипом <<<P000069>,
для каждого типа <<<P000070> и <<P000071>.
\Item{\SrnTypeVariableReflexivityA}} Продвигаемая переменная
Тип <<<P000072> Это подтип <<<P000073>.
<<<P000954>{<<P000955>> Правильно повышен Переменная А}
Тип <<<P000074> является подтипом <<<P000075> если
<<P000076> является подтипом как <<<P000077>, так и <<<P000078>.
<<<P000956>{<<<P000957>> Правильное будущее Или А.Тип <<<P000079> является подтипом <<<P000595> если
<<P000081> Это подтип <<<P000596>.
\Item{<<P000959>>>}} Правильное будущее Или Б.
Тип <<<P000083> является подтипом <<<P000597> если
<<P000085> Это подтип <<<P000086>.
\Item{\SrnLeftPromotedVariable}} Продвигаемая левая переменная B}
Тип <<<P000087> является подтипом <<<P000088> если
<<P000089> является подтипом <<<P000090>.
\Item{\SrnLeftVariableBound}} Левая переменная граница
Переменная типа <<<P000091> является подтипом типа <<<P000092> если
Связь <<<P000093>
(как указано в текущей среде <<<P000094>)
Это подтип <<<P000095>.
\Item{\SrnRightFunction}} Правильная функция
Каждый тип функции является подтипом типа <<<P000966>.
\Item{\SrnPositionalFunctionType}} Тип позиционной функции
Тип функции <<<P000096> с позиционными дополнительными параметрами
является подтипом
Другой тип функции <<<P000097> с опциональными параметрами
Если первый имеет самое
количество требуемых параметров, равное последнему,
Последний имеет по крайней мере
общее количество параметров, равное первому;
Тип возврата <<<P000098> является подтипом подтипа <<<P000099>;
и каждого типа параметров <<<P000100> <<<P000969>{супертип}
соответствующий тип параметров <<<P000101>, если таковой имеется.
Обратите внимание, что отношение к типам функций без дополнительных параметров
и отношения между типами функций без дополнительных параметров,
Покрывается разрешением <<<P000102> соответственно <<<P000103>.
Для каждого подтипа отношения, рассматриваемого в этом правиле,
Формальные параметры типа <<<P000104> и <<<P000105> должны быть приняты во внимание
(как это отражено в использовании расширенной среды <<<P000106>).
Мы можем предполагать без потери общности.
что имена переменных типов попарно идентичны,
Мы считаем, что типы общих функций эквивалентны в
последовательное переименование
(<<P000523>).
Короче говоря, «во время доказательства мы переименуем их по мере необходимости».
Наконец, обратите внимание, что взаимосвязь между негенерическими типами функций
Покрывается разрешением <<<P000107>.
\Item{\SrnNamedFunctionType}} Названы функции типа
Тип функции <<<P000108> с названными факультативными параметрами является подтипом
Другой тип функции <<<P000109> с названными дополнительными параметрами
если они имеют одинаковое количество требуемых параметров,
и набор имен названных параметров для последнего является подмножеством
Для тех, кто раньше
Тип возврата <<<P000110> является подтипом подтипа <<<P000111>;
и каждый параметр типа <<<P000112> <<<P000972>{супертип}
соответствующий тип параметров <<<P000113>, если таковой имеется.
Обратите внимание, что отношение к типам функций без дополнительных параметров
и отношения между типами функций без дополнительных параметров,
Покрывается разрешением <<<P000114> соответственно <<P000115>,
и что последний случай идентичен правилу, полученному из
Правило ~\SrnPositionalFunctionType
в отношении подтипирования между типами функций без дополнительных параметров.
Как правило, <<<P000974>,
Мы можем принять без потери общности.
Имена переменных типов попарно идентичны.
Аналогично, негенерические функции охватываются разрешением <<<P000116>.
\Item{<<P000976>>>}} Класс Ковариантность
Параметризованный тип на основе общего класса <<<P000117> является подтипом
параметризованный тип на основе того же класса <<<P000118> если
Каждый действительный тип аргумента первого является подтипом
соответствующего фактического типа аргумента последнего.
Это правило может иметь значение <<<P000119> и охватывают также необычный класс,
Но это избыточно, потому что это уже покрыто
Правило <<<P000977>.
\Item{\SrnSuperinterface}} СуперинтерфейсРассмотрим случай, когда <<<P000120> и <<<P000121> Во-первых,
параметризованный тип на основе негенерического класса <<<P000122> является подтипом
параметризованный тип на основе другого негенерического класса <<<P000123> если
<<P000124> Это прямой суперинтерфейс <<<P000125>.
Когда <<<P000126> или <<<P000127>, это правило описывает отношение подтипа
который включает один или несколько общих классов,
В этом случае мы должны дать названия формальным параметрам типа <<<P000128>,
и указать, как они используются в спецификации суперинтерфейса
На основе <<P000129>.
С этими частями на месте, мы можем указать отношения подтипа.
Существует между двумя параметризованными типами на основе <<<P000130> и <<<P000131>.
%
% % TODO(Eernst): Обратите внимание, что спецификация того, как передавать аргументы типа
% % в <<<P000524> Это неправильно, и это должно быть
% полностью переписаны для интеграции новой конструкции смесителя.
Случай, когда суперкласс представляет собой приложение для смешивания, рассматривается через
эквивалентность с объявлением
обычный (возможно, общий) суперкласс
(<<P000525>),
Это означает, что может быть несколько шагов подтипа от
a заданное объявление класса классу, указанному в <<<P000980>{} пункт.
<<<P000500>%
?


<<<P000981>{дополнительные концепции субтипирования}
<<P000838>

<<<P000785>%
<<P000132> <<<P000554> <<<P000133> в данной среде <<<P000134>,
написанный <<P000982>{S}{T}
если {f} <<P000983> T.S.

<<<P000786>%
Тип <<P000135>
<<P000555>
Тип <<<P000136> в среде <<<P000137>,
написанный <<P000984>{S}{T}
Если {f} <<P000985>{S}{T} или <<P000986> T.S.
В этом случае мы говорим, что типы <<<P000138> и <<<P000139> есть
<<P000556>.

<<<P000987>{%>
Это правило может удивить читателей, привыкших к обычной проверке типа.
Намерение <<<P000988> отношение
Это не означает, что задание гарантирует динамический успех.
Вместо этого он направлен только на флаги.
Это почти наверняка будет ошибочным.
Не исключая заданий, которые могут работать.

Например, назначение объекта статического типа <<<P000598>
переменной со статическим типом <<P000599>,
Несмотря на то, что это не гарантировано,
Может быть хорошо, если значение времени выполнения оказывается строкой.

Статический анализатор или компилятор
может поддерживать более строгие статические проверки в качестве опции.
?


<<<P000989>{Функциональные типы}
<<P000839>

<<<P000787>%
<<<P000990>{Функциональные типы}{тип!функция}
Выпускаются в двух вариантах:

<<P000486>
<<P000991>
Типы функций, которые имеют только позиционные параметры.
Они имеют общую форму
<<<P000992>{T}.
<<P000993>
Виды функций с названными параметрами.
Они имеют общую форму
<<<P000994>{T}.
<<P000501>

<<<P000995>{%>
Обратите внимание, что необычный случай охватывается наличием <<<P000140>,
в этом случае объявления о параметрах типа опускаются
(P000526).
Случай без факультативных параметров покрывается наличием <<<P000141>;
Обратите внимание, что все правила, касающиеся типов функций двух типов
совпадает в данном случае.%
?

<<P000788>%
Два типа функций считаются равными при последовательном переименовании типа.
Параметры могут сделать их идентичными.

<<<P000996>{%>
Обычный способ сказать, что мы не различаем
функции, которые являются альфа-эквивалентными.
Правило субтипирования ниже означает, что мы можем предположить, что
Подходящее переименование уже состоялось.
В тех случаях, когда это невозможно
Поскольку количество типовых параметров в двух типах различается
Или же границы разные,
Подтипа не существует.%
?

<<<P000789>%
Объект функции всегда является экземпляром некоторого класса <<<P000142> который реализует
класс <<<P000997>< (<<P000527>),
и который имеет метод, названный <<P000998>,сигнатурой которого является тип функции <<<P000143> себя.
<<<P000999>{%>
Следовательно, все типы функций являются подтипами <<<P001000>
(<<P000528>)%
?


\subsection{Тип \FUNCTION}}
<<P000840>

<<<P000790>%
Встроенный класс <<<P001003>{} является супертипом всех типов функций.
(<<P000529>).
Невозможно расширить, реализовать или смешать в классе <<<P001004>.

<<<P000791>%
Если заявление класса или приложение для смешивания имеет <<<P001005>{} в качестве суперкласса,
Вместо этого он использует P000600 в качестве суперкласса.

<<<P000792>%
Если класс или смешанная декларация реализуется <<<P001006>, это не имеет никакого эффекта.
Это как если бы P001007 Он был удален из <<P000601> пункт
(и если это единственный реализованный интерфейс, весь пункт удаляется).
Полученный класс или интерфейс микширования
не имеет <<<P001008>{} в качестве суперинтерфейса.

<<<P000793>%
Если приложение для смешивания смешивает <<<P001009>{} на суперкласс, следует за
Нормальные правила применения миксина, но так как результат этого миксинга
Приложение эквивалентно классу с <<<P000602>, и
Это положение не имеет эффекта, полученный класс также не
осуществлять <<P001010>.
<<<P001011>{%>
<<<P001012>{} класс не объявляет конкретных членов,
Таким образом, приложение микширования создает подкласс
Суперкласс без новых членов и без новых интерфейсов.
?

<<<P001013>{%>
После использования \FUNCTION{} Таким образом, это не имеет никакого эффекта, это было бы
разумно полностью запретить его, как мы расширяем, реализуем или
смешивание в таких типах, как <<P000603> или <<P000604>.
Для обратной совместимости с программами Dart 1
Синтаксис разрешается сохранять, даже если он не имеет эффекта.
Инструменты могут предупреждать пользователей о том, что их код не имеет эффекта.
?


\subsection{Тип <<<P001016>>
<<P000841>

<<<P000794>%
Тип <<<P001017>{} является статичным типом, который является супертипом всех других типов.
Как будто <<P000605>,
отличается от других видов тем, что статический анализ
Предполагает, что каждый член имеет соответствующий член
с подписью, допускающей данный доступ.

<<<P001018>{%>
Например,
когда приемник в обычном способе вызова имеет тип <<<P001019>,
Можно использовать любое название метода,
с любым количеством аргументов или без них,
с любым количеством аргументов,
И всякие именованные доводы,
любого типа,
Без ошибок.
Обратите внимание, что вызов все равно вызовет ошибку времени компиляции.
Если есть ошибка в одном или нескольких аргументах или других терминах.
?

<<<P000795>%
Предполагается, что процентный вывод имел место, поэтому тип не был выведен.
Если аннотация статического типа не была предоставлена,
Типовая система рассматривает декларации как имеющие тип <<<P001020>.
% % TODO(Eernst): Изменение при добавлении спецификации instantiate-to-bound.
Если используется общий тип, но аргументы типа не приводятся,
Аргументы типа по умолчанию вводят <<<P001021>.

<<<P001022>{%>
% % TODO(Eernst): Изменить при добавлении спецификации мгновенного подключения.
Это означает, что при наличии общей декларации
<<P000606>,
где <<P000147> является формальной декларацией параметров типа <<<P000148>,
Тип <<<P000149> является эквивалентным

<<<P001023>
<<<P000607>>%
?

<<P000796>%
Встроенная декларация типа <<P000608>,
который объявлен в библиотеке <<P000609>,
обозначает \DYNAMIC{} Тип.
Когда имя \DYNAMIC{} Экспорт: <<P000610> оценивается
как выражение,
Для сравнения: <<P000611> Объект, представляющий собой \DYNAMIC{} тип,
Даже если <<<P001027>{} не является классом.

<<<P001028>{%>
<<<P000612>> объект должен быть равен соответствующему <<<P000613>
<<<P000614> и <<P001029>>> По словам оператора \lit{==================================================================================================================
(<<P000530>)%
?

<<P000797>%
Для повышения точности статических типов
Член получает доступ к приемнику типа \DYNAMIC{} который относится к
Декларации встроенного класса <<P000615>дается статический тип, соответствующий этим декларациям
Всякий раз, когда это происходит, это звучит хорошо.

<<P000487>
<<P001032>
<<<P000152> быть выражением формы <<<P000616>,
который не следует за аргументом,
где статический тип <<<P000154> <<<P001033>,
<<P000769> имя геттера, объявленное в <<<P000617>;
если тип возврата <<<P000618> <<<P000155> затем
Статический тип <<<P000156> является <<<P000157>.
<<<P001034>{%>
Например, <<P000619> имеет тип <<P000620>
<<P000621> имеет тип <<<P000622>,%
?

<<<P001035>
<<<P000158> быть выражением формы <<<P000623>,
который не следует за аргументом,
где статический тип <<<P000160> <<<P001036>,
<<P000770> это название метода, объявленного в <<<P000624>
сигнатура метода которого имеет тип <<<P000161>
<<<P001037>{(который является функциональным типом)}.
Статический тип <<<P000162> В этом случае <<P000163>.
<<<P001038>{%>
Например, <<P000625> имеет тип <<<P000626>.%
?

<<P001039>
<<<P000165> быть выражением, которое имеет форму
\code{$d$.\id(\metavar{аргументы})}
или форму \code{$d$.\id<\metavar{typeArguments}>(\metavar{arguments})},
где статический тип <<<P000168> <<<P001047>,
<<P000771> имя геттера, объявленное в <<<P000627> Тип возврата <<P000169>,
\metavar{аргументы} происходит от \synt{аргументы} и
\metavar{typeArguments} являются производными от \synt{typeArguments}, если они присутствуют.
Статический анализ обрабатывается <<<P000170> как функция выражения призыв
где объект статического типа <<<P000171> Применяется к данной аргументной части.
<<<p001052>{%>
Это всегда ошибка времени компиляции.
Например, <<P000628> Ошибка времени компиляции,
Потому что он проверяется как
Выражение функции вызова, где
сущность статического типа <<P000629> это
Призвано. Обратите внимание, что это действительно может быть успешным: Преобладающая реализация
<<<P000630> Можно ли вернуть пример
динамический тип которого является подтипом
<<<P000631>> Имеет метод <<<P001053>{}.
Мы решили сделать это ошибкой, потому что это может быть ошибкой.
Особенно в таких случаях, как <<P000632>
где разработчик мог забыть, что <<P000633> Геттер.%
?

<<<P001054>
<<<P000174> быть выражением формы <<<P001055>{<<<P000175>. \id(\metavar{аргументы})}
где статический тип <<<P000176> <<<P001058>, <<<<P001059>{аргументы}
фактический список аргументов, полученный из <<<P001060>{аргументов},
<<P000772> Это название метода, объявленного в <<<P000634>
чей метод подписи имеет тип <<P000177>.
Если число позиционных фактических аргументов в <<<P001061>
меньше числа требуемых позиционных аргументов <<<P000178>
или больше числа позиционных аргументов в <<<P000179>,
или если <<<P001062>{аргументы} включает любые названные аргументы
с именем, которое не указано в <<<P000180>,
Тип <<<P000181> <<<P001063>.
В противном случае тип <<<P000182> Тип возврата в <<<P000183>.
<<<P001064>{%>
<<<P000635> имеет тип <<<P001065>< тогда как
<<P000636> имеет тип <<P000637>.
Обратите внимание, что призывы, которые «не подходят» статически
Заявления не являются ошибками, они просто получают тип возврата. <<<P001066>>%
?

<<P001067>
<<<P000186> быть выражением формы
<<<P001068><<<P000187>. \id<\metavar{typeArguments}>(<<P001071>>>{arguments})} где
Статический тип <<<P000188> - <<<P001072>,
<<<P001073>{typeArguments} - список актуальных
Аргументы типа, полученные из <<<P001074>{typeArguments} и<<<P001075>{аргументы} - это фактический список аргументов, полученный из \synt{аргументы}.
Ошибка времени компиляции <<<P000773> Это имя
необычный метод, указанный в <<P000638>.
<<<P001077>{%>
Общие методы не указаны в <<<P000639>.
Поэтому мы не утверждаем, что должно быть
Статически необходимое количество аргументов фактического типа и
Они должны соответствовать границам.
В противном случае это был бы последовательный подход.
Потому что призыв гарантированно потерпит неудачу, когда любой из этих
нарушены требования,
Но обобщения этого механизма должны включать такие правила.
?

<<P001078>
Например, вызов метода <<<P000189> (в том числе призывы геттеров,
сеттеры и операторы), где приемник имеет статический тип <<<<P001079>{}
<<P000190> не соответствует ни одному из вышеперечисленных случаев статического типа <<<P000191> это
<<P001080>.
Когда выполняется выражение, полученное из <<<P001081>
геттер или метод вызова, который соответствует одному из приведенных выше случаев,
Применяется соответствующий статический анализ и ошибки компиляции времени.
<<<P001082>{%>
Например, <<P000640> Ошибка.%
?
<<P000502>

<<<P001083>{%>
Обратите внимание, что только очень немногие формы вызова метода экземпляра с
приемник типа \DYNAMIC{} Это может быть ошибка времени компиляции.
Конечно, некоторые выражения, такие как <<P000641> Синтаксические ошибки
Хотя они также могут считаться «призывами»,
Субэкспрессии проверяются отдельно.
Любой конкретный аргумент может быть ошибкой времени компиляции.
Но почти любая форма списка аргументов может быть обработана через
<<P000642>,
Любой аргумент может быть принят, потому что любой
Формальный параметр в основной декларации может иметь свой тип
Аннотация была изменена на <<<P000643>.
Это естественное следствие принципа
это
a Получатель <<<P001085>{} допускает почти все вызовы метода экземпляра.
Несколько случаев, когда метод инстанций вызывает
Получатель типа <<<P001086>{} Это ошибка
Либо гарантированно проваливаются во время бега,
или они очень, очень вероятно, будут ошибки разработчика.
?


<<<P001087> Тип будущего или
<<P000842>

<<P000798>%
Встроенная декларация типа <<P000644>,
который вывозится библиотекой <<P000645>,
определяет общий тип с одним параметром типа (<<P000531>).
Тип <<<P000646> Это неклассовый тип
Для всех, кто имеет регулярные ограничения <<<P000194>.

<<<P001088>{%>
Отношения подтипа с участием <<<P000647> указывается в других
(<<P000532>).
Обратите внимание, однако, что они влекут за собой определенные полезные свойства:

<<P000488>
<<<P001089> <<P000195>
<<<P000196>T<<<P000197>.
<<<p001090>> <<P000198>
<<P000199> T<<<P000200> T<<<P000201>.
<<<P001091>[<<P000202>>>]
Если <<<P000203> и <<<P000204> T$>} <: S$, затем $\code{FutureOr<$ T<<<P000207>.
<<P000503>

То есть <<<P000648> является в некотором смысле
Союз <<<P000208> и соответствующего будущего типа.
Последний пункт гарантирует, что
<<<P000649> <:<<P000650>,
А также <<<P000651> является ковариантным по своему типу параметру,
Как и типы классов:
Если <<<P000210> <:<<P000211>, то <<P000652> <: <<P000653>
?

<<P000799>%
Если аргументы типа перешли на <<<P000654> будет иметь ошибки компиляции времени
если применяется к обычному общему классу с одним параметром типа,
Ошибки времени компиляции выдаются для <<<P000655>.
Оригинальное название: P000656 как выражение
обозначает <<<P000657> объект, представляющий тип <<P000658>.

<<<P001092>{%>
<<<P000659> Тип представляет собой случай, когда объект может быть
Или один из экземпляров типа <<P000215>
или тип <<P000660>.
Такие случаи происходят естественным образом в асинхронном коде.
Доступной альтернативой будет использование верхнего типа (например, <<P001093>).<<<P000661> позволяет использовать некоторые инструменты для
Более точный анализ типов.%
?

<<<P000800>%
Тип <<P000662> имеет интерфейс, идентичный этому
<<P000663>.
<<<P001094>{%>
То есть только члены, которые <<<P000664>> можно ли ссылаться
на объекте со статическим типом <<<P000665>,%
?

<<<P001095>{%>
Мы допускаем только призывы членов, которые унаследованы от
Обычный супертип как <<<P000219>, так и <<<P000666>.
В большинстве случаев единственным распространенным супертипом является <<P000667>.
Исключения, такие как <<P000668>
Что такое P000669 как обычный супертип,
малочисленны и практически не полезны,
Поэтому сейчас мы выбираем только призывы к
члены, унаследованные от <<<P000670>.%
?

<<<P000801>%
Определяем вспомогательную функцию
<<<P001096><<<P001097> T}}{futureOrBase(t)@\emph{futureOrBase}$(T)$}
следующим образом:

<<P000489>
<<P001099> Если <<<P000222> <<<P000671> Для некоторых <<P000224>
<<<P000225>.
<<<P001100> Иначе <<P000226>.
<<P000504>


<<<P001101> Тип Пустота.
<<P000843>

<<<P000802>%
Специальный тип \VOID{} Используется для обозначения
Значение выражения бессмысленно и должно быть отброшено.

<<<P001103>{%>
Типичным случаем является тип \VOID{} Используется как тип возврата
Функция, которая «ничего не возвращает».
Технически всегда будет объект <<<P001105>{some}
который возвращается
(<<P000533>).
Но очень важно иметь функцию
чьей единственной целью является создание побочных эффектов,
так, что <<<P001106>{любое} использование возвращаемого объекта
Это было бы неправильно.
%
Это не значит, что есть что-то неправильное
Возвращенный объект как таковой.
Это может быть любой объект.
Но разработчик, выбравший тип возврата <<<P001107>
Это указывает на то, что это недоразумение
придать этому объекту какой-либо смысл,
или использовать его в любых целях.%
?

<<<P001108>{%>
Тип <<<P001109> является высшим типом
(<<P000534>),
\VOID{{} и <<P000672> являются подтипами друг друга
(<<P000535>),
Это означает, что любой объект может быть
значение выражения типа <<<P001111>.
%
Следовательно, любой экземпляр типа <<<P000673> который соответствует типу \VOID{}
должны сравниваться равными (по оператору <<<P001113>{==} <<<P000536>)
Для любого случая <<P000674> который соответствует типу <<<P000675>
(P000537).
Не гарантируется, что <<P000676> Оценивается как правда.
На самом деле, не рекомендуется, чтобы реализации стремились достичь этого.
Потому что может быть более важным обеспечить, чтобы диагностические сообщения
(включая следы стека и динамические сообщения об ошибках)
Хранить достаточно информации, чтобы использовать слово «пустота» при обращении к типам
которые указаны как таковые в исходном коде.%
?

<<<P000803>%
В поддержку этого понятия
значение выражения со статическим типом \VOID{} должны быть отброшены,
Такие объекты могут использоваться только в конкретных ситуациях:
Возникновение выражения типа \VOID{} является ошибкой времени компиляции
Если это не разрешено в соответствии с одним из следующих правил.

<<P000490>
<<<P001116>
В <<<P001117>{выражение <<P000677>, <<P000228> Может иметь тип <<<P001118>.
<<<P001119> Значение <<<P000229> отбрасывается.
<<P001120>
При инициализации и инкрементных выражениях петли,
\code{<<P001122>>>{) (<<P000230>; <<P000231>; <<P000232> {\ldots}}
<<P000233> может иметь тип <<<P001124>,
и каждое из выражений в списке выражений <<<P000234> может иметь тип <<P001125>.
<<<P001126> Значения <<<P000235> и <<<P000236> отбрасываются.
<<P001127>
В типе литья <<<P000678>, <<P000239> может иметь тип <<P001128>.
<<<P001129>{%>
Таким образом, разработчики получают возможность <<<P001130>>на значениях с статическим типом <<<P001131>.
Это означает, что такие ценности не должны быть отброшены,
Но они могут быть использованы только тогда, когда это необходимо.
Указывается явно.%
?
<<P001132>
В скобчатом выражении <<<P000679>, <<P000241> Может иметь тип <<<P001133>.
<<<P001134>{%>
Обратите внимание, что <<<P000680> Сам по себе тип <<P001135>,
Это означает, что оно должно происходить в определенном контексте.
где это не ошибка, чтобы иметь его.%
?
<<P001136>
В условном выражении <<<P000681>,
<<P000246> и <<P000247> может иметь тип <<P001137>.
<<<P001138>{%>
Статический тип условного выражения тогда <<<P001139>,
Даже если одна из ветвей имеет другой тип,
Это означает, что условное выражение должно произойти снова.
В каком-то контексте, где это не ошибка, чтобы иметь его.
?
<<P001140>
В нулевом слитном выражении <<<P000682>,
<<P000250> Может иметь тип <<<P001141>.
<<<P001142>{%>
Статический тип нулевого коалесцирующего выражения тогда <<<P001143>,
Что, в свою очередь, ограничивает, где это может произойти.
?
<<P001144>
В выражении формы <<<P000683>, <<P000252> может иметь тип <<P001145>.
<<<P001146>{%>
Это правило было принято, потому что это было существенное изменение.
Превратить эту ситуацию в ошибку
В то время, когда была изменена обработка <<<P001147>.
Инструменты могут дать подсказку в таких случаях.%
?
<<P001148>
<<<P001149>{%>
В ответном заявлении <<<P000684>,
<<P000254> может иметь тип <<<P001150>{} в ряде ситуаций
(<<P000538>)%
?
<<P001151>
<<<P001152>{%>
В теле функции стрелы <<<P000685>,
Возвращенное выражение <<P000256> может иметь тип <<<P001153><
в ряде ситуаций
(<<P000539>)%
?
<<P001154>
Начальное выражение для переменной типа \VOID{}
может иметь тип <<P001156>.
<<<P001157> Использование этой переменной ограничено.
<<P001158>
Фактическое выражение аргумента, соответствующее формальному параметру
Статически известный тип аннотации <<<P001159>
может иметь тип <<P001160>.
<<<P001161>{%>
Использование этого параметра в теле вызывающего
Статически ожидается, что они будут ограничены наличием типа <<P001162>.
Смотрите дискуссию о звуке ниже
(<<P000540>)%
?
Это правило также применяется к операторам <<<P000686> и <<<P000687>.
<<<P001163>{%>
Например, <<<P000257> и $e_2$ может иметь тип <<<P001164><
В выражении формы <<P000688>
когда параметры статически известного оператора <<<P000689>
Оба имеют тип <<<P001165>{.%}
?
Наконец, это правило относится и к сеттерам.
<<<P001166>{%>
Например, с выражением формы <<<P000690>
Который обозначает призывник,
Статически известный параметр типа <<<P001167>,
<<P000264> может иметь тип \VOID{.%}
?
<<P001169>
В выражении формы <<<P000691>
где <<P000267> является <<<P001170>{назначаемый выражение, обозначающее локальную переменную
<<<P001171>{(который также может быть формальным параметром)}
тип <<P001172>,
<<P000268> может иметь тип <<P001173>.
<<<P001174>{%>
Использование этой переменной ограничено, поскольку она имеет тип <<<P001175>.
Смотрите дискуссию о звуке ниже
(<<P000541>)%
?
<<P001176>
<<<P000269> быть выражением, заканчивающимся на <<<P001177>
в форме <<P000692>,
где <<P000273> имеет форму

<<P001178>
\syntax{(<cascadeSelector> <argumentPart>*)
(<assignableSelector> <argumentPart>*)*)

<<P001180>
<<<P000274> Имеет форму <<<P001181>{выражение БезКаскада}.

Если <<<P000275> является <<<P001182>{назначенный выборщик}
форма <<P000693> или <<P000694>где статический тип идентификатора <<<P000774> <<<P001183>,
<<P000276> может иметь тип <<P001184>.
В противном случае, если <<<P000277> <<<P001185>{назначаемый выборщик} формы
<<P000695> где статический тип
Первый формальный параметр в статически известной декларации
оператор <<<P000696> <<<P001186>,
<<P000279> может иметь тип <<P001187>.
Кроме того, если статический тип второго формального параметра <<<P001188>,
<<P000280> Может иметь тип <<<P001189>.
<<P000505>

<<<P000804>%
Наконец, мы должны рассмотреть ситуации, связанные с неявным использованием
Объект, статический тип которого может быть <<<P001190>.
%
Это ошибка времени компиляции для заявления for-in, чтобы иметь итератор.
выражение типа <<<P000281> <<<P001191>{Итератор\VOID{}>}
Наиболее специфичная инстанциация <<<P000697>
Это суперинтерфейс <<<P000282>, если только
Переменная итерации имеет тип <<P001193>.
%
Это ошибка времени компиляции для асинхронного заявления for-in.
Иметь потоковое выражение типа <<P000283>
<<<P001194>{Stream\VOID{}>} является наиболее специфическим
Инстанциация <<<P000698> Это суперинтерфейс <<<P000284>,
Если переменная итерации не имеет типа <<P001196>.

<<<P001197> Вот несколько примеров:

<<P000000>

<<<P001198>{%>
Однако в примерах, которые не являются ошибками
Использование <<P000699> В петле тело ограничено,
Потому что имеет тип <<<P001199>%
?


<<<P001200> Звуковая пустота
<<P000844>

<<<P000805>%
Ограничения на использование объекта, полученные в результате оценки
выражение со статическим типом \VOID{}
не строго соблюдаются.

<<<P001202>{%>
Использование «пустого значения» не является вопросом здравости, то есть
нет инвариантов, необходимых для правильного выполнения программы Дарта
могут быть нарушены из-за такого использования.
?

<<<P001203>{%>
Можно сказать, что тип <<<P001204> используется
Помогать разработчикам поддерживать определенную дисциплину
О том, что некоторые объекты не предназначены для использования.
%
В силу того, что принуждение не является необходимым,
и из-за обработки <<<P001206>{} в более ранних версиях Дарта,
Язык использует подход <<<P001207>{лучшие усилия} для обеспечения
что значение выражения типа \VOID{}
не используется.%
?

<<<<P001209>{%>
На самом деле, есть много способов в дополнение к типу литья.
где разработчик может получить доступ к такому объекту: %
?

<<<P000001>

<<<<P001210><%
В (1) переменная <<<P000700> тип <<<<P001211 >>{}
общая функция <<<P000701>,
что допускается, так как фактический тип аргумента \VOID{}
и разрешается передавать фактический аргумент типа <<<P001213><<
Формальный параметр того же типа.
%
Тем не менее, никакой специальной обработки не дается, когда выражение имеет тип.
которая является или содержит переменную типа, значение которой может быть <<<P001214>,
Поэтому нам разрешено вернуться <<<P000702> в теле <<<P000703>,
Это означает, что мы косвенно получаем доступ к стоимости.
выражение типа <<<P001215>, под статичным типом <<P000704>.

В пункте 2 мы косвенно получаем доступ к значению
Переменная <<<P000705> с типом <<<P001216>,
потому что мы используем задание, чтобы получить доступ к экземпляру <<P000706>
который был создан с аргументом типа <<<P001217> по типу
<<P000707>.
Обратите внимание, что \code{A<Object>} и \code A<\VOID{}> являются подтипами друг друга,
То есть они эквивалентны по правилам подтипа,
Таким образом, ни статические, ни динамические проверки не будут работать.

В пункте 3 мы косвенно получаем доступ к стоимости
Переменная <<<P000709> с типом <<<P001220><
Статический тип <<<P000710>,
потому что статически известный способ подписи <<<P000711>
имеет тип параметра <<<P001221>,Но фактическая реализация <<P000712> на который ссылается
является оверрайдом, тип параметра которого <<<P000713>,
что допускается, потому что <<<P000714> и \VOID{} являются двумя верхними типами.
?

<<<P001223>{%>
Очевидно, что поддержка для предотвращения разработчиков от использования объектов
полученных из выражений типа \VOID{} Далеко не звук,
В том смысле, что есть много способов обойти правило
Такой объект следует выбросить.

Однако мы решили сосредоточиться на простом использовании первого порядка.
(где выражение имеет тип <<<P001225>, и используется значение),
И мы оставили дела более высокого порядка в значительной степени неконтролируемыми,
использование дополнительных инструментов, таких как linters, для проведения анализа;
Покрывает косвенные потоки данных.

Конечно, можно было бы определить разумные правила.
Таким образом, значение выражения типа \VOID{}
будет гарантированно выброшен после некоторых переводов.
от одной переменной или параметра к следующей, все с типом <<<P001227>,
явно или в качестве значения параметра типа.
В частности, мы можем потребовать, чтобы переопределение метода
Никогда не отменяйте тип возврата <<<P000715> Тип возврата <<<P001228>,
Типы параметров в противоположном направлении;
параметризованные типы с аргументом типа \VOID{} не может быть назначен
переменных, где аргумент соответствующего типа является чем-либо иным, чем
<<<P001230> и т.д.\ etc.

Но это было бы совершенно непрактично.
В частности, необходимость предотвращения большого количества переменных типа.
из когда-либо имевших значение <<<P001231>,
или предотвращение использования определенных значений, тип которых является такой переменной типа,
или в котором содержится такая переменная типа,
Это сильно ограничило бы большую часть всего кода Дарта.

Поэтому мы решили помочь разработчикам поддерживать эту самонавязанную дисциплину.
в простых и прямых случаях,
и оставить его для специальных рассуждений или отдельных инструментов для обеспечения
Косвенные случаи рассматриваются так же тщательно, как это необходимо на практике.
?


<<<P001232> Параметризованные типы
<<P000845>

% TODO (Eernst): Возможно, мы хотели бы добавить новую концепцию применения
% общая функция для аргументов фактического типа (возможно, это дополнительный вид аргументов).
% "параметризованный тип", но он отличается от общего класса, поскольку
% мы можем иметь динамические вызовы общей функции. Но это делает
% не возникает как самостоятельная организация до того, как мы введем общие отрывы
% (<<P000716>), или если мы позволим ему возникнуть на неявной основе
% на вывод. Эта новая концепция, вероятно, должна быть добавлена в этот раздел.

<<<P000806>%
<<<P001233>{параметризированный тип} — синтаксическая конструкция.
где применяется наименование декларации общего типа
Список фактических аргументов.
<<<P001234>{общая инстанциация} - операция, при которой
Общий тип применяется к фактическим аргументам типа.

<<<P001235>{%>
Параметризованный тип — это синтаксическая концепция, которая соответствует
Семантическая концепция родовой инстанциации.
При использовании первого мы часто оставляем последний неявным.
?

<<<P000807>%
<<<P000285> быть параметризованным <<P000717>.

<<<P000808>%
Ошибка времени компиляции <<<P000288> не является общим типом,
<<<P000289> является универсальным типом,
но количество формальных параметров типа в декларации <<<P000290> Это не P000291.
В противном случае пусть
<<P000292>
быть формальными параметрами типа <<<P000293>, и
<<P000294>
быть соответствующими верхними границами, используя <<<P001236> когда не объявлялось никаких обязательств.

<<<P000809>%
<<<P000295> <<<P000557> Если {f}
<<P000296> Для одного или нескольких <<P000297>,
<<<P000298> Это не очень хорошо (P000542).

<<<P000810>%
Ошибка времени компиляции <<<P000299> Он ограничен.

<<<P000811>%
<<<P000300> оценивается следующим образом.
<<<P000301> быть результатом оценки <<<P000302>, для <<<P000303>.<<P000304> Затем оценивает дженерикальную инстанциацию
где <<P000305> Применяется к <<<P000306>.

<<<P000812>%
<<<P000307> быть параметризованным типом формы
<<P000718>
Предположим, что <<P000310> Он не деформирован и не деформирован.
Если <<P000311> является статическим типом элемента <<<P000312> <<<P000313>,
Статический тип элемента <<<P000314> Выражение типа
<<P000719>
это
<<P000317>,
где <<P000318> Формальные параметры типа <<<P000319>.

% % TODO(Eernst): Это место, где мы можем указать, что каждый из
% % аргументы типа получателя <<P000720> должен быть
% % заменено нижним типом («Null», на данный момент) в местах расположения члена
Тип %%, где это происходит противоположно. Например, "c.f" должен иметь
% % статический тип «пустая функция» когда "c" имеет статический тип "C<T>" для любого
% % 'T', и мы имеем 'класс C<X> { void Function(X) f; }'.


<<<P001237>{Актуальные типы}
<<P000846>

% Примечание (Eernst): Аргументы фактического типа в этом разделе являются динамическими объектами.
% не синтаксис (используется понятие «фактический тип» и «фактическая связь»)
% определяют динамическую семантику, включая динамические ошибки. Мы используем <<P000322>
% для обозначения аргументов такого типа, как и все те места, где
Используется понятие %, а не <<P000323> который часто используется для обозначения
% синтаксис аргумента реального типа.
%
% Дело в том, что <<P000324> никогда не будет содержать формальный параметр типа, поэтому
% не нужно беспокоиться о необходимости повторять замену, которая находит
% фактической границы. Например, нет проблем с определением фактической границы.
% для <<P000721> в обращении к <<<P001238>{void foo< X расширяет C<X>>() <<<P001489>...<<<P001490>>
%, даже если это рекурсивное обращение по форме <<P000722> в
% тела.

<<<P000813>%
Пусть <<<P000325> Термин, полученный из <<<P001239>{тип}.
Пусть <<<P001240>{X}{1}{s} будут формальными параметрами типа в области применения
в месте, где <<P000326> происходит.
В контексте, где фактический тип аргументов соответствует
<<<P001241> X}{1}{s} являются <<<P001242>{t}{1}{}
<<<P000558> <<P000327> тогда
<<P000328>.

<<<P000814>%
<<<P000329> быть декларацией с именем <<P000330> и введите аннотацию <<<P000331>.
<<<P000559> <<<P000332> и <<P000333> В данном контексте является
Тогда фактическое значение <<P000334> в этом контексте.

<<<P001243>{%>
В необычный случай, когда <<<P000335>,
фактический тип равен заявленному типу,
В том смысле, что мы используем простые термины, такие как <<P000723>, для обозначения обоих.
Обратите внимание, что <<<P001244>{X}{1}{s} может быть объявлено несколькими объектами, например,
одна или несколько прилагаемых общих функций и прилагаемый общий класс.%
?

<<<P000815>%
<<<P001245><<<P000336> <<<P001246>< <<<P000337>> Формальное объявление параметров типа.
<<<P000560> <<<P000338> В данном контексте является
Фактическое значение <<P000339> в этом контексте.

<<<P001247>{%>
Обратите внимание, что даже если <<<P000340> Может возникнуть в <<<P000341>
не происходит в фактическом значении <<P000342>,
Поскольку ни один тип не имеет фактического значения, которое включает переменную типа.
?


<<<P001248> Наименее верхние границы
<<P000847>

% TODO (Eernst): Этот раздел был обновлен для выполнения общих функций.
%, но никаких других изменений не произошло. Следовательно, мы все еще
% необходимо обновить этот раздел, чтобы использовать правила Dart 2 для LUB.

<<<P000816>%
В некоторых случаях это расходится?
Дано два интерфейса <<P000343> и <<P000344>,
<<<P000345> быть набором суперинтерфейсов <<<P000346>,
<<<P000347> Набор суперинтерфейсов <<<P000348>
И пусть <<P000349>.
Кроме того,
<<<P000350>> Для любого конечного числа <<P000351>
где <<P000352> Число шагов на самом длинном пути наследования
От <<P000353> до <<P000724>.%TODO(lrn): Укажите: «путь наследования» — это путь в графе суперинтерфейса.
<<<P000354> Самое большое число, такое, что <<P000355> Кардинальность одна,
Он должен существовать, потому что <<P000356> является <<P000357>.
Наименьшая верхняя граница <<<P000358> и <<<P000359> является единственным элементом <<<P000360>.

<<<P000817>%
Наименьшая верхняя граница <<<P001249> и любой тип <<<P000361> является <<<P001250>.
Наименьшая верхняя граница <<<P001251> и любой тип <<<P000362> является <<<P001252>.
Наименьшая верхняя граница <<P000363> Любой тип <<<P000364> является <<<P000365>.
<<<P000366> быть переменной типа с верхней границей <<<P000367>.
Наименьшая верхняя граница <<P000368> Тип <<P000369> это
наименьшая верхняя граница <<<P000370> и <<<P000371>.

<<<P000818>%
Операция с наименьшей верхней границей является коммутативной и идемпотентной.
Но она не является ассоциативной.

% Типы функций

% % TODO(Eernst): Этот раздел должен использовать новый синтаксис для типов функций.
%% (<<<P001253>{} и т.д.); эти обновления сделаны CL 84908.

<<<P000819>%
Наименее верхняя граница типа функции и типа интерфейса <<<P000372> это
наименьшая верхняя граница <<<P001254>{} и <<<P000373>.
<<<P000374> и <<<P000375> быть функциональными типами.
Если <<P000376> и <<P000377> отличаются по количеству требуемых параметров,
Тогда наименьшая верхняя граница <<P000378> и <<P000379> составляет <<<P001255>.
В противном случае:

<<P000491>
<<P001256> Если

<<P001257>
<<P000725> и

<<<P001258>
<<P000726>

<<P001259>
где <<P000390> наименьшая верхняя граница <<<P000391> и <<<P000392> это

<<P001260>
<<P000727>

<<P001261>
где <<P000397> является наименьшей верхней границей <<<P000398> и <<<P000399>.
<<P001262> Если

<<P001263>
<<P000728>,

<<P001264>
<<<P001265><<<P000405> $X_1\ B_1, \ldots,\ X_s\ B_s$><<<P000407><<<P001491> <<<P001266>< <<<P001492>> <<P000408>>

<<P001267>
наименьшая верхняя граница <<<P000409> и <<<P000410> это

<<P001268>
<<P000729>

<<P001269>
где <<P000414> является наименьшей верхней границей <<<P000415> и <<<P000416>.
<<<P001270> Если

<<P001271>
<<<P001272><<<P000417> $X_1\ B_1, \ldots,\ X_s\ B_s$>$T_1, \ldots,\ T_r,\ $\{$T_{r+1}\ p_{r+1}, \ldots,\ T_f\ p_f$\} <<P000421>,

<<P001273>
<<<P001274><<<P000422> $X_1\ B_1, \ldots,\ X_s\ B_s$>$S_1, \ldots,\ S_r,\ $\{$S_{r+1}\ q_{r+1}, \ldots,\ S_g\ q_g$\} <<<P000426>>

Тогда пусть
<<P000427>
и пусть <<<P000428> быть наименьшей верхней границей типов <<<P000429> в <<<P000430> и
<<P000431>.
Тогда наименьшая верхняя граница <<<P000432> и <<<P000433> это

<<P001275>
<<<P001276><<<<P000434>> $L_1, \ldots,\ L_r,\ $\{$X_m\ x_m, \ldots,\ X_n\ x_n$\} <<P000437>

где <<P000438> является наименьшей верхней границей <<<P000439> и <<<P000440>
<<P000506>

<<<P001277>{%>
Обратите внимание, что необычный случай охватывается использованием <<<P000441>,
в этом случае объявления параметров типа опущены (<<P000543>).%
?


<<<P001278> Ссылка.
<<P000848>


<<<P001279>{Лексические правила}
<<P000849>

<<<P000820>%
Исходный текст Dart представлен в виде последовательности точек кода Unicode.
Эта последовательность сначала преобразуется в последовательность токенов
в соответствии с лексическими правилами, приведенными в данной спецификации.
В любой момент процесса токенизации,
Признается самый длинный токен.


<<<p001280>> Зарезервированные слова
<<P000850>

<<<P000821>%
<<<P000561> Может использоваться только в синтаксических положениях.
определяется грамматикой.
В частности, ошибка времени компиляции возникает при использовании зарезервированного слова.
где ожидается идентификатор.

<<<P001281>{%>
Обратите внимание, что зарезервированные слова встречаются жирными и нецитируемыми в правилах грамматики.
(например, <<<P001282>{})
Даже если последовательное обозначение будет использовать цитаты
(например, <<<P001283>{assert}).
Это злоупотребление происходит потому, что мы считаем, что
Это делает правила грамматики более читаемыми.
?

<<P000492><RESERVED<<<P001499> Слово>:= \ASSERT{} | \BREAK{} | \CASE{} | \CATCH{} !
\CLASS{} | \CONST{}
\alt\hspace{-3mm} \CONTINUE{} | \DEFAULT{} | \DO{} | \ELSE{} | \ENUM{} !
\EXTENDS{} | \FALSE{} | \FINAL{} | \FINALLY{} | \FOR{}
\alt\hspace{-3mm} \IF{} | \IN{} | \IS{} | \NEW{} | \NULL{} | \RETHROW{} !
\RETURN{} | \SUPER{} | \SWITCH{} | \THIS{}
\alt\hspace{-3mm} \TRUE{} | <<P001318>>>{} | <<P001319>>>{} | \VOID{} | \WHILE{} | <<P001322>>>{}
<<P000507>

<<<P000822>%
В грамматике должно происходить правило для зарезервированных слов выше.
перед правилом для <<<P001323> BUILT<<<P001500> В <<<P001501> Идентификатор
(P000544).

<<<P001324>{%>
Это гарантирует, что \synt{IDENTIFIER} и <<P001326>>>{IDENTIFIER<<P001502>>> No<<<P001503> ДОЛЛАР. Не надо.
Они выводят любые зарезервированные слова и не выводят никаких встроенных идентификаторов.
?


<<<P001327> Комментарии
<<P000851>

<<<P000823>%
\IndexCustom{Комментарии}{комментарий}
Это разделы текста программы, которые используются для документации.

<<P000493>
<<<<P001504> LINE<<<P001505> Комментарий> ::= <<<P001329><
<<<P001330> (<LINE\_BREAK>)* (перенаправлено с «P001507») Дырнуть?

<MULTI<<<P001508> LINE<<<P001509> Комментарий> ::= \gnewline{}
"/*" (<MULTI<<P001510>>>) LINE<<<P001511> Комментарий> | \gtilde{} '*/''** "
<<P000508>

<<<P000824>%
Dart поддерживает как однострочные, так и многострочные комментарии.
<<<P000562> начинается с токена <<P000730>.
Все, что между <<P000731> и конец строки
Компилятор Dart должен быть проигнорирован
Если комментарий не является комментарием к документации.

<<<P000825>%
<<<P000563> начинается с токена <<P000732>
И кончается токеном <<P000733>.
Все, что между <<<P000734>* и <<<P000735>/
Компилятор Dart должен быть проигнорирован
Если комментарий не является комментарием к документации.
Комментарии могут гнездиться.

<<<P000826>%
\IndexCustom{Комментарии к документации}{комментарии к документации}
Комментарии, которые начинаются с токенов <<P000736> или <<P000737>.
Комментарии к документации предназначены для обработки
Инструмент, который производит человеческую читаемую документацию.

<<<P000827>%
Текущий объем комментария к документации, непосредственно предшествующий
Объявление класса <<<P000442> Это
<<<P001334>{Объем тела}{Сфера действия! для органа по заявлениям
<<P000443>.

<<<P000828>%
Текущий объем комментария к документации, непосредственно предшествующий
Объявление неперенаправляющего генеративного конструктора <<<P000444>
с параметрами инициализации - формальный параметр инициализации <<<P000445>
(<<P000545>).

<<<P000829>%
В противном случае текущая сфера применения комментария к документации, непосредственно предшествующего
Объявление функции <<<P000446> Формальный диапазон параметров <<<P000447>
(P000546).


<<<P001335> Прецедент оператора.
<<P000852>

<<<P000830>%
Преимущество оператора косвенно определяется грамматикой.

<<<P001336>{%>
Следующая ненормативная таблица может быть полезной
<<P001337>

<<<P000002>%
?


<<P001338> Оригинальное название: Algorithmic Subtyping
<<P000853>

% Номер правил подтипа
\newcommand{\AppSrnReflexivity>>>{<<<P001341>> 1_{\scriptsize\mbox{algo}}}}
\newcommand{\AppSrnTypeVariableReflexivityB}{<<<P001346>. 1)
\newcommand{\AppSrnTypeVariableReflexivityC}{<<<P001349>. 2.
\newcommand{<<P001351>>>}{<<P001352>>>.3.
\newcommand{\AppSrnRightFutureOrC}{\SrnRightFutureOrB. 1)
\newcommand{\AppSrnRightFutureOrD}{\SrnRightFutureOrB. 2.

<<<P000494>[h!]
\def\VSP{\vspace{3 мм)
\def\ExtraVSP{\vspace{1mm)
\def\Axiom#1#2#3#4{<<P001367>>>{<<P001368>>>[#1]{}{<<P001369>>>{#3}{#4}}}<<< P001370>>>\def\Rule#1#2#3#4#5#6{\centerline{<<P001374>>>[#1]{<<P001375>>>{#3}{#4}}{<<P001376>>>{#5}{#6}}}<<<< P001377>>>
\def<<<P001379>#1#2#3#4#5#6#7#8
<<<<<<<<<<<<<<<<<<<<<<<<<<<<<<<<<<<<<<<<<<<<<<<<<<<<<<<<<<<<<<<<<<<<<<<<<<<<<<<<<<<<<<<<<<<<<<<<<<<<<<<<<<<<<<<<<<<<<<<<<<<<<<<<<<<<<<<<<<<<<<<<<<<<<<<<<<<<<<<<<<<<<<<<<<<<<<<<<<<<<<<<<<<<<<<<<<<<<<<<<<<<<<<<<<<<<<<<<<<<<<<<<<<<<<<<<<<<<<<<<<<<<<<<<<<<<<<< P001385>>>
<<<P001386><<<P001387>#1#2#3#4#5
\centerline{<<P001389>>>[#1]{#3}{<<P001390>>>{#4}{#5}}}<<< P001391>>>
<<<P001392><<<P001393>#1#2#3#4<<<P001394><<<P001395> [#1]{#3}{#4}}\VSP]
%
\begin{minipage}[c]{0.49<<P001397>>>}
\RuleRaw{\AppSrnReflexivity>>> Рефлексивность {S\mbox{ Не композитный. S}{S}
\Rule{\AppSrnTypeVariableReflexivityC>>> Тип переменной рефлексивности B}{X}{T}{X}{X <<<P001512> T
\Rule{\AppSrnRightFutureOrC>>> Правильное будущее Или C}{\Delta(X)}{<<P000738>>>}{X}{<<P000739>>>}
<<P000509>
<<P000496>>>[c]{0.49<<P001406>>>}
\Axiom{\AppSrnTypeVariableReflexivityB>>> Переменная рефлексивность типа {X}{X}
<\Rule{\AppSrnTypeVariableReflexivityD>>> Переменная рефлексивность типа C {X} <<<P001513> S}{T}{X <<<P001514> S {\displaystyle X} <<P001515> T
<\Rule{\AppSrnRightFutureOrD>>> Правильное будущее Или D}{S}{<<P000740>{X} <<<P001516> S}{\code{FutureOr<<<<P000451>>>>}
<<P000510>
%
<<<P001413> Правила алгоритмического подтипа.
Правила \SrnTop--\SrnSuperinterface{} неизменны и поэтому здесь опущены.
<<P000549>
<<P000511>

<<<P000831>%
Текст этого приложения не является частью спецификации языка Дарта.
Тем не менее, мы по-прежнему используем запись, где точная информация
использует стиль, связанный с нормативным текстом в спецификации (этот стиль);
<<<P001416>{в то время как примеры и объяснения используют стиль комментариев (как этот)}.

<<<P000832>%
В настоящем добавлении представлен вариант правил подтипа, приведенных
На рисунке ~<<<P000547> на странице ~<\pageref{фиг:правила подтипа}.

<<<P001418>{%>
Правила докажут тот же набор отношений подтипа,
Но приведенные здесь правила показывают, что существует эффективная реализация.
Это будет определять, будет ли <<<P001419>{S}{T} держать,
для любых типов <<P000452> и <<P000453>.
Легко видеть, что алгоритмические правила докажут самое большее.
одинаковые отношения подтипа,
Все приведенные здесь правила могут быть доказаны.
С помощью правил на рисунке ~<<<P000548>.
Также относительно легко нарисовать доказательства.
Алгоритмические правила могут доказать, по крайней мере, те же отношения подтипа.
также при соблюдении следующих ограничений по порядку и прекращению.%
?

<<P000833>%
Единственным измененным правилом является число ~<<<P001420>.
Изменяется на <<<P001421>.
Это только меняет применимость правила:
Это правило используется только для типов, которые не являются атомными.
<<<P001422>{атомный} Тип {type!atomic}
Это тип, который не является переменной типа.
не является продвигаемой переменной типа,
не функциональный тип,
Не параметризованный тип.

<<<P001423>{%>
Иными словами, правило <<<P001424>{} используется для
Особые типы, такие как <<<P001425>, <<<P001426> и <<<P001427>,
Используется для негенеральных классов.
Но он не используется для любого типа операций.
Для этого требуется более одного сравнения, чтобы определить,
Это то же самое, что и любой другой тип.
%
Дело в том, что оставшиеся правила заставят
в любом случае, при необходимости,
Таким образом, мы можем просто опустить начальный структурный переход.
которые могут предпринять много шагов только для того, чтобы сообщить об этих двух терминах большого типа.
Не совсем идентичны.%
?

<<<P000834>%
Правила упорядочены по их номерам правил:
Приведенное здесь правило пронумеровано <<<P000454> вставляется сразу после правила <<<P000455>,
Далее следует правило <<<P000456>, и так далее.
Далее следует правило, число которого <<<P000457>.
<<<P001428>{%>
Так что приказ
<<<P001429>, \SrnTop--\SrnTypeVariableReflexivityA,
<<<P001432>, <<<P001433>,
<<<P001434>,
<<<P001435>> и так далее.%
?

<<<P000835>%
Теперь мы определим процедуру, которая используется для определения того,
<<<P001436>{S}{T} держит,для некоторых конкретных типов <<<P000458> и <<P000459>:
Выберите первое правило <<P000460> чьи синтаксические ограничения удовлетворены
по данным типам <<P000461> и <<P000462>,
и далее показать, что его помещения держатся.
Если это так, мы заканчиваем и заключаем, что отношения подтипа имеют место.
В противном случае мы прекращаем и завершаем
что отношения подтипа не поддерживаются,
За исключением случаев, когда <<P000463> это
<<P001437>, <<<P001438>,
<<<P001439>, или <<P001440>.
<<<P001441>{%>
В частности, для исходного запроса \SubtypeStd Сюда!
Мы не отказываемся от использования правила, которое
более высокое число правил, чем у <<<P000464>,
Но мы можем попробовать все
Правила с <<<P000742> вправо.%
?

<<<P001443>{%>
Помимо того, что вся сложность субтипирования
потенциально возникает каждый раз, когда проверяется, имеет ли предпосылка
Проверки, применяемые для каждого правила, связаны с количеством работы.
которые являются постоянными для всех правил, кроме следующих:
Во-первых, группа правил
<<<P001444>, <<<P001445>,
<<<P001446> и <<P001447>>>
Это может привести к обратному ходу.
Далее, правила \SrnPositionalFunctionType-\SrnCovariance{}
требует работы, пропорциональной размерам <<<P000466> и <<<P000467>,
из-за количества помещений, которые необходимо проверить.
Наконец, правило ~\SrnSuperinterface требует работы
пропорционально размеру <<<P000468>,
Это также может повлечь за собой затраты на поиск всего набора данных.
прямые и косвенные суперинтерфейсы подтипа кандидата <<P000469>,
До тех пор, пока не будет показана соответствующая предпосылка для одного из них,
Если есть.

Применяются дополнительные оптимизации.
Например,
Сразу можно сделать вывод, что отношения подтипа не
когда мы собираемся проверить правило ~<<<P001451>
Если <<P000470> Это переменная типа или тип функции.
Для некоторых других типов, например,
- переменная продвигаемого типа,
<<P000743>, <<<P001452>, <<<P001453>>
<<<P000744> для любого <<<P000472> или <<<P001454>,
Известно, что это никогда не произойдет как <<P000473> для правила ~<<<P001455>,
Что означает, что этот, казалось бы, дорогой шаг
может быть ограничено до некоторой степени.%
?


<<<P001456> Добавление: Целые варианты осуществления
<<P000854>

<<<P001457>{%>
<<<P000745> Тип представляет целые числа.
Спецификация написана с целыми числами дополнения 64-битных двух в качестве
предполагаемой реализации. Но когда Дарт компилируется в JavaScript,
Использование <<P000746> Вместо этого он будет использовать тип номера JavaScript.
и соответствующих операций JavaScript,
За исключением битовых операций, как описано ниже.

Это приводит к ряду различий:

<<P000497>
<<<P001458> <<P000474>
Действительные значения JavaScript <<P000747> IEEE-754 64-битный
Число с плавающей точкой без дробной части. Это включает
Положительный и отрицательный <<<P000564>, который может быть достигнут
переполнение (целое деление на ноль по-прежнему является динамической ошибкой).
В противном случае действительные целые буквы (включая любой ведущий знак минус)
Представляет собой недействительный JavaScript <<P000748> Значения составляют время компиляции
Ошибки при компиляции на JavaScript. Операции на целых числах могут
теряют точность, потому что операнды и результат представлены
64-битные числа с плавающей запятой, ограниченные 53-значными
кусочки.
<<<P001459> <<P000475>
JavaScript <<P000749> Также используются экземпляры <<<P000750> и
целочисленное значение <<<P000751> Примеры также реализуют <<P000752>.
<<P000753> и <<P000754> Классы остаются отдельными подклассами
класса <<<P000755>, но <<<P001460>{инстанции} любого класса, который
представляют собой целое действие, как если бы они на самом деле были
Общий подкласс, использующий оба <<P000756> и <<P000757>.
Фракционные числа реализуют только <<P000758>.
<<<P001461>[<<P000476>]Битовые операции на целых числах (\lit{\&}, <<<P001463>{ |}} <<<P001464><<<P001518>>,
<<<P001465><<<P001466>> \lit{<<P001468>>>} и <<P001469>>>{<<P001470>>>} Все усеченные
операнды к целым числам комплемента 32-битных двух, выполняют 32-битные
операции на них, и полученные 32 бита интерпретируются как
32-битное неподписанное целое число. Например, <<P000759> оценивать
4294967294 (перенаправлено с «P000760»). Правый
оператор смены, \lit{\gtgt}, выполняет подписанное правое смещение
32 бита, когда исходное число отрицательное, и неподписанное право
смещение, когда исходное число не является отрицательным. Оба вида сдвига
Читать далее <<P000477> в положение <<P000478>, <<<P000479>; но подписано
сдвиг оставляет бит 31 неизменным, а неподписанный сдвиг записывает 0
Как раз 31. Например:
<<<P000761>>
<<P000762>.
В этом примере мы отмечаем, что <<<P000763> и <<<P000764>
получить представление дополнения 32-битных двух <<<P000765>,
Но они имеют разные значения для бита знака IEEE 754.
<<<P001473>[<<P000480>>>]
<<P000766> Способ не может различать значения <<<P000481> и
<<P000482>, и он не может распознать любой <<P000565> значение как идентичное
Для себя. Для эффективности: <<P000767> Операция использует
JavaScript Оператор P000768.
<<<P000512>%
?

<<<P001474>
